\begin{table}[htbp]
\centering
\caption{Effect of MEI Yes-Vote Share on Foreign Population Share}
\label{tab:main_did}
\begin{tabular}{lcccc}
\hline\hline
 & (1) & (2) & (3) & (4) \\
\hline
MEI Yes-Share (std.) $\times$ Post & 0.0020*** & 0.0021*** & 0.0019*** & 0.0013** \\
 & (0.0005) & (0.0005) & (0.0007) & (0.0006) \\
\hline
Municipality FE & Yes & Yes & Yes & Yes \\
Year FE & Yes & Yes & --- & --- \\
Pre-FS $\times$ Year & & Yes & Yes & \\
Canton $\times$ Year FE & & & Yes & \\
Language $\times$ Year FE & & & & Yes \\
Within $R^2$ & 0.004 & 0.012 & 0.007 & 0.001 \\
Observations & 29,372 & 29,372 & 29,372 & 29,372 \\
\hline\hline
\end{tabular}
\begin{tablenotes}[flushleft]\footnotesize
\item \textit{Notes:} Standard errors clustered at the municipality level in parentheses. The dependent variable is the foreign population share (mean = 0.165). MEI Yes-Share is standardized (mean 0, SD 1). A one-unit increase corresponds to an 11.3 percentage point higher anti-immigration vote share. Specification (2) adds pre-2014 foreign share interacted with year dummies. Specification (3) adds canton $\times$ year fixed effects. Specification (4) replaces canton FE with language region $\times$ year FE. $^{*}p<0.1$; $^{**}p<0.05$; $^{***}p<0.01$.
\end{tablenotes}
\end{table}

